% Draft for the B-owned supervised methods subsection. Integrate after verifying
% the final frozen protocol and the team-wide LaTeX macros and bibliography.
\paragraph{BiLSTM--CRF (M3).}
We concatenate a frozen 100-dimensional GloVe word embedding with a
character-level CNN representation for each word. The character encoder uses
learned character embeddings, a width-three convolution, and max pooling;
unlike the lowercase GloVe lookup, it retains the original character case.
A bidirectional LSTM contextualizes these representations, and a linear-chain
CRF models the tag sequence. This architecture tests a learned contextual
representation while retaining sequence-level decoding.

\paragraph{Fine-tuned BERT (M4).}
We fine-tune \texttt{bert-base-cased} for token classification. Words may be
split into subwords; supervision and prediction use the first subword of each
word, while noninitial subwords and special tokens do not contribute to the
classification loss. BIO tags that start with an invalid inside tag are
repaired before extracting entities. This model tests the effect of a large
pretrained contextual encoder.

\paragraph{Span classifier.}
We enumerate candidate word spans of at most 14 words. A fine-tuned BERT
encoder provides start and end representations, which are concatenated with a
learned span-length embedding and classified into an entity type or the
non-entity class. Training retains positive spans and samples negative spans;
inference greedily resolves overlapping positive predictions. This model tests
entity-level classification without BIO sequence decoding. It only predicts
flat, nonoverlapping entities and cannot recover a span longer than its
candidate-length limit.

% Add the final search space, chosen settings, seeds, data versions, metrics,
% confidence intervals, and citations from the frozen experiment manifest.
